\documentclass[10pt,a4paper]{amsart}
\usepackage[margin=19mm]{geometry}
\usepackage{amsmath,amssymb,mathtools}
\usepackage[scr=rsfs,cal=euler]{mathalfa}
\usepackage{xfrac}
\usepackage[unicode,hidelinks,bookmarks=false]{hyperref}
\newcommand{\ie}{i.e.\ }
\newcommand{\cf}{cf.\ }
\newcommand{\resp}{respectively\ }
\newcommand{\Subsection}{\subsection}
\providecommand{\nobreakdash}{\nobreak\hbox{-}}
\setlength{\emergencystretch}{2em}
\multlinegap0pt
\usepackage{amssymb}
\usepackage{dsfont}
\usepackage{enumerate}
\usepackage{bm}
\usepackage{xcolor}
\newtheorem{thm}{Theorem}
\newtheorem{rem}{Remark}
\newcommand{\eps}{{\varepsilon}}
\title{Enhanced stability and asymptotic limits to the non-isentropic compressible fluid-particle interaction model with thermal effects}
\author{Fucai Li, Jinkai Ni, Zhouping Xin}
\date{Source: arXiv:2607.17115v1 (CC BY 4.0)}
\begin{document}
\maketitle

\noindent\emph{Source and licence.} Enhanced stability and asymptotic limits to the non-isentropic compressible fluid-particle interaction model with thermal effects. Version arXiv:2607.17115v1 (CC BY 4.0), distributed by arXiv under the Creative Commons Attribution 4.0 International licence, http://creativecommons.org/licenses/by/4.0/, as recorded in the arXiv licence metadata for this exact version and shown on the abstract page https://arxiv.org/abs/2607.17115v1. Full licence notice: \texttt{licenses/arxiv-2607.17115-CC-BY-4.0.txt}.\par\smallskip

\noindent\textit{Editorial scope.} Counted unit: the article's thm T1.4 (source0.tex lines 910--953). The article's complete original proof of it is delivered with it (source0.tex lines 3726--3879, one \texttt{proof} environment of 154 lines, which the article places away from the statement). No mathematics is written, repaired or re-derived: the statement and the proof are the article's own lines, in the article's own notation, and no retained line is substituted or re-flowed.  Retained uncounted as the source-local context the proof needs: lines 651--664; lines 666--670; lines 776--778; lines 1357--1361; lines 1756--1766.  Nothing was omitted from any retained span, so the package carries no omission record.  No retained line contains a citation, so the wrapper carries no bibliography and no bibliography entry.  No retained line contains a figure environment, an image-inclusion command or drawing code, so no image is delivered and the package declares no asset dependency.  Editorial formatting only, in addition: an AMS \texttt{amsart} preamble that restates the article's own declarations and macros used by the retained lines, namely the environments \texttt{thm}, \texttt{rem} on separate counters, together with the standard packages those lines need.  The retained environments print in this document as Theorem 1, Remark 1 in order of appearance, whereas the article prints the same items with the numbering of its own full document.  The complete native source of the article as distributed in the arXiv e-print for this exact version is bundled unchanged as \texttt{sources/205535/source0.tex}.
\begin{equation}\label{Q2}
\left\{\begin{aligned}
&\partial_t \rho + u\cdot\nabla\rho+(1+\rho){\rm div} u=0, \\
&\partial_t u +u \cdot \nabla u+\frac{1+\theta}{1+\rho} \nabla 
\rho+\nabla \theta=\frac{b-u-au}{1+\rho}, \\
&\partial_t \theta +u \cdot \nabla \theta+\theta  {\rm div} u
+ {\rm div} u-\sqrt{3}(\sqrt{2} \omega-\sqrt{3} \theta)  \\
&\quad =\frac{ |u|^2-2 u \cdot b+a|u|^2-3 a \theta }{1+\rho}  
 -\frac{\sqrt{3}\rho(\sqrt{2} \omega-\sqrt{3} \theta)}{1+\rho},\\
&\partial_t f +v \cdot \nabla f+u \cdot \nabla_v f-\frac{1}{2} u \cdot v f-
u \cdot v \sqrt{M}- (|v|^2-3 ) \sqrt{M} \theta  \\
&\quad =\mathcal{L} f+\frac{\theta}{\sqrt{M}} \Delta_v(\sqrt{M} f),\\
\end{aligned}\right.
\end{equation} 

\begin{align}\label{Q2-2}
(\rho ,u ,\theta ,f )|_{t=0}=&\,(\rho _0(x),u _0(x),\theta _0(x),f _0(x,v))\nonumber\\
:=&\,\Big( \varrho_0 (x)-1,u _0(x), {\vartheta} _0(x)-1,\frac{F _0(x,v)-M}{\sqrt{M}}\Big),
 \quad   x \in \mathbb{R}^3, \, v\in\mathbb{R}^3.  \quad
\end{align}

\begin{align}\label{b-1}
\mathcal{E}_0:=\|(\rho _0,u _0,\theta _0 )\|_{H^4}^2+\|f_0 \|_{H_{x,v}^4}^2\leq \eps_1,
\end{align}

\begin{align} 
-\langle \mathcal{L}g, g\rangle &\geq \bar\lambda|\{\mathbf{I}-\mathbf{P}_0\}g|_{\nu}^{2}, \label{G2.6a}  \\  
-\langle \mathcal{L}\{\mathbf{I}-\mathbf{P}\}g,g\rangle &\geq \bar\lambda|\{\mathbf{I}-\mathbf{P}\}g|_{\nu}^{2}, \label{G2.6b}  \\  
-\langle \mathcal{L}g, g\rangle &\geq \bar\lambda |\{\mathbf{I}-\mathbf{P}\}g|_{\nu}^{2} + |b|^{2} + 2|\omega|^2.\label{G2.6}  
\end{align}

\begin{equation}\label{G3.21}
\left\{\begin{aligned}
&\partial_t a+{\rm div} b=0,\\
& \partial_t b_i+\partial_i a+\frac{2}{\sqrt{6}} \partial_i \omega+\sum_{j=1}^3 \partial_{j} \Gamma_{i, j}(\{\mathbf{I}-\mathbf{P}\} f)=-b_i+u_i+u_i a, \\
& \partial_t \omega+\sqrt{2}(\sqrt{2} \omega-\sqrt{3} \theta)-\sqrt{6} a \theta+\frac{2}{\sqrt{6}} {\rm div}  b-\frac{2}{\sqrt{6}} u \cdot b    +\sum_{i=1}^3 \partial_{i} \Upsilon_i(\{\mathbf{I}-\mathbf{P}\} f)=0, \\
& \partial_j b_i+\partial_i b_j- (u_i b_j+u_j b_i )  -\frac{2}{\sqrt{6}} \delta_{i j}\bigg(\frac{2}{\sqrt{6}} {\rm div}   b-\frac{2}{\sqrt{6}} u \cdot b+\sum_{i=1}^3 \partial_{ i} \Upsilon_i(\{\mathbf{I}-\mathbf{P}\} f)\bigg) \\
&\quad =-\partial_t \Gamma_{i, j}(\{\mathbf{I}-\mathbf{P}\} f)+\Gamma_{i, j}(\mathfrak{l}+\mathfrak{r}+\mathfrak{s}), \\
& \frac{5}{3} (\partial_i \omega-\omega u_i-\sqrt{6} \theta b_i )-\frac{2}{\sqrt{6}} \sum_{j=1}^3 \partial_{j} \Gamma_{i, j}(\{\mathbf{I}-\mathbf{P}\} f)   \\
&\quad=-\partial_t \Upsilon_i(\{\mathbf{I}-\mathbf{P}\} f)+\Upsilon_i(\mathfrak{l}+\mathfrak{r}+\mathfrak{s}).
\end{aligned}\right.
\end{equation}

\begin{thm}\label{T1.4}
%Let $\Omega=\mathbb{R}^3$ and 
% $\kappa=0$.
 Suppose that the assumption
\eqref{b-1} holds and $\|(\rho_0,u_0,\theta_0)\|_{L^1}+\|f_0\|_{ {Z}_1}$ is bounded, then
the solution $(\rho,u,\theta,f)$ to the Cauchy problem \eqref{Q2}--\eqref{Q2-2} satisfying
 \begin{equation}\label{d1} 
\left\{
\begin{aligned}
 \|\nabla^k(\rho, u, \theta)(t)\|_{L^2 } +\|\nabla^kf(t)\|_{{L_v^2(L^{2})}}
&\leq C_3(1+t)^{-\frac{3}{4}-\frac{k}{2}},\quad k=0,1, \\
\|\nabla^k(\rho, u, \theta)(t)\|_{L^2 } +\|\nabla^kf(t)\|_{{L_v^2(L^{2})}}
&\leq C_3(1+t)^{-\frac{5}{4}},\quad\,\,\,\,\,\,\, k=2,3,4
\end{aligned}
\right.
\end{equation}
and
 \begin{equation} \label{d2}
\left\{
\begin{aligned}
\| (\rho, u, \theta)(t)\|_{L^p}+\|  f(t)\|_{L_v^2(L^p)}
 &\leq C_3(1+t)^{-\frac{3}{2}(1-\frac{1}{p})},\,\,\, 2\leq p\leq 6, \\  
\| (\rho, u, \theta)(t)\|_{L^p}+\|  f(t)\|_{L_v^2(L^p)}
 &\leq C_3(1+t)^{-\frac{5}{4} },\,\,\,\quad\quad\,  6\leq p\leq \infty. 
\end{aligned}
\right.
\end{equation}
Furthermore, we have 
\begin{align}\label{NJKd3}
 \|  (b-u)(t)\|_{L^2}+\|  (\sqrt{2}\omega-\sqrt{3}\theta)(t)\|_{L^2}
 +\|  \{\mathbf{I}-\mathbf{P}\}f(t)\|_{L_v^2(L^2)}  \leq C_3(1+t)^{-\frac{5}{4} }.
\end{align}
where    $C_3>0$ is a constant independent of    $t$. 

\begin{rem}
 We observe that the decay rate of   $L^2$-norm for the dissipative 
 quantities $ b - u $ and $ \sqrt{2}\omega - \sqrt{3}\theta $ is
  $ (1 + t)^{-\frac{5}{4}} $, which is  $\frac{1}{2}$-order faster 
  than the decay rate of the solution itself. This  new phenomenon is induced by the coupling effects between 
 the particles and the fluid and is constituted   first time for non-isentropic compressible fluid-particle models. 
\end{rem} 
 

\end{thm}

\begin{proof}[Proof of Theorem \ref{T1.4} \emph{(}continued\emph{)}]
To end the proof of  Theorem \ref{T1.4},  we first need to investigate  
the  time-decay rates of classical solutions in $L^p$ space with  $p\in [2,+\infty]$. In fact, 
based on the decay rates obtained in   \eqref{d1}, we have 
\begin{align*}
\|(\rho,u,\theta)\|_{L^2}\lesssim \,& (1+t)^{-\frac{3}{4}}, \qquad
 \|(\rho,u,\theta)\|_{L^6}\lesssim   \|\nabla(\rho,u,\theta)\|_{L^2}\lesssim (1+t)^{-\frac{5}{4}},   \\
 \|f\|_{L_v^2(L^{2})}\lesssim \,&  (1+t)^{-\frac{3}{4}},\qquad\quad\, 
 \| f\|_{L_v^2(L^{6})}\lesssim    \|\nabla f\|_{L_v^2(L^{2})}\lesssim (1+t)^{-\frac{5}{4}}.      
\end{align*}

For $p\in [2,6]$, using the interpolation inequality, it holds that
\begin{align*}
 \|(\rho,u,\theta)\|_{L^p}\lesssim \,
 &  \|(\rho,u,\theta)\|_{L^2}^{\zeta}\|(\rho,u,\theta)\|_{L^6}^{{1-\zeta}}
 \lesssim  (1+t)^{-\frac{3}{2}(1-\frac{1}{p})} , \\
 \|f\|_{L_v^2(L^{p})}\lesssim \,&  
 \|f\|_{L_v^2(L^{2})}^{\zeta}\|f\|_{L_v^2(L^{6})}^{{1-\zeta}}\lesssim  (1+t)^{-\frac{3}{2}(1-\frac{1}{p})},
\end{align*}
where $\zeta= ({6-p})/ {2p} \in [0,1]$.

Using   Gagliardo-Nirenberg's inequality, we derive  that
\begin{align*} 
\|(\rho,u,\theta)\|_{L^{\infty}} \lesssim \,& \|\nabla(\rho,u,\theta)\|_{H^{1}} 
\lesssim (1+t)^{-\frac{5}{4}}, \\ 
\|f\|_{L_v^2(L^{\infty})} \lesssim \,& \|\nabla f\|_{L_v^2(H^{1})} 
\lesssim (1+t)^{-\frac{5}{4}}. 
\end{align*}  

For $p \in [6,\infty]$, by applying the interpolation inequality, it follows that  
\begin{align*} 
\|(\rho,u,\theta)\|_{L^p} &\lesssim \|(\rho,u,\theta)\|_{L^6}^{\zeta'} 
\|(\rho,u,\theta)\|_{L^{\infty}}^{1-\zeta'} \lesssim (1+t)^{-\frac{5}{4}}, \\ 
\|f\|_{L_v^2(L ^{p})} &\lesssim \|f\|_{L_v^2(L ^{6})}^{\zeta'} 
\|f\|_{L_v^2(L ^{\infty})}^{1-\zeta'} \lesssim (1+t)^{-\frac{5}{4}}, 
\end{align*}  
where $\zeta' = {6}/{p} \in [0,1]$. This guarantees that \eqref{d2} holds.

 \medskip 
 Finally, we prove the remaining inequality \eqref{NJKd3}. 
 It follows  from \eqref{Q2}$_2$, \eqref{Q2}$_3$ and 
\eqref{G3.21}$_1$ \linebreak  –\eqref{G3.21}$_3$ that
\begin{align}\label{NJKG3.158}
 \partial_t (b-u)+2(b-u) 
=\,&  -\nabla a-\frac{2}{\sqrt{6}}\nabla\omega-{\rm div}
\big\langle\{ v\otimes v-{\rm Id}\} \sqrt{M},\{\mathbf{I}-\mathbf{P}\}f               
    \big\rangle+2au\nonumber\\
&+\nabla \rho+\nabla\theta  +u\cdot\nabla u
+h(\rho)\theta\nabla\rho +g(\rho)  \big(  \nabla\rho +(b-u)+au\big)\nonumber\\
\equiv:\,& \mathcal{Q}_1, 
\end{align}
and
\begin{align}\label{NJKG3.159}
&\partial_t(\sqrt{2}\omega-\sqrt{3}\theta)+5(\sqrt{2}\omega-\sqrt{3}\theta)\nonumber\\ 
=\,& \sqrt{2}\Big(\sqrt{6}a\theta-\frac{2}{\sqrt{6}}{\rm div}b
+\frac{2}{\sqrt{6}}u\cdot b-{\rm div} \Big\langle \frac{v(|v|^2-3)}{\sqrt{6}}\sqrt{M}, 
\{\mathbf{I}-\mathbf{P}\}f   \Big\rangle   \Big)\nonumber\\
&+\sqrt{3}{\rm div}u+\sqrt{3}\theta{\rm div}u
+\sqrt{3} u\cdot\nabla\theta-\sqrt{3} u\cdot b-3\sqrt{3}a\theta\nonumber\\
&-\sqrt{3} h(\rho) \big( -u\cdot(b-u) +a|u|^2 
+\sqrt{3}(\sqrt{2}\omega-\sqrt{3}\theta)-b\cdot u-3a\theta      \big)  \nonumber\\
\equiv:\,& \mathcal{Q}_2.
\end{align}
The equations \eqref{NJKG3.158} and \eqref{NJKG3.159} can be represented as % follows:
\begin{align}\label{NJKG3.160}
b-u=&\, e^{-2t}(b_0-u_0)+\int_{0}^t e^{-2(t-\tau)} \mathcal{Q}_1{\rm d}\tau,    
\end{align}
and
\begin{align}\label{NJKG3.161}
\sqrt{2}\omega-\sqrt{3}\theta=&\, e^{-5t}(\sqrt{2}\omega_0-\sqrt{3}\theta_0)
+\int_{0}^t e^{-5(t-\tau)} \mathcal{Q}_2{\rm d}\tau.   
\end{align} 
Taking the $L^2$-norm of both sides of \eqref{NJKG3.160} and \eqref{NJKG3.161}, and applying 
the decay estimates in \eqref{d2}, we obtain   
\begin{align}\label{NJKG3.162}
\|(b-u)(t)\|_{L^2}\lesssim&\, e^{-2t} \|b_0-u_0\|_{L^2}+\int_0^t e^{-2(t-\tau)}
\big(\|\nabla(a,\rho,\theta,\omega)\|_{L^2}+\|\nabla f\|_{L_{x,v}^2}\big)  {\rm d}\tau\nonumber\\
&+\int_0^t e^{-2t}\big(\|au\|_{L^2}+\|u\cdot\nabla u\|_{L^2}
+\|\theta\nabla\rho\|_{L^2}+\|\rho\|_{L^\infty}\|b-u\|_{L^2}  \big){\rm d}\tau \nonumber\\
\lesssim&\, e^{-2t}+\int_0^t e^{-2(t-\tau)}(1+\tau)^{-\frac{5}{4}}{\rm d}\tau 
+\Big(\int_0^t e^{-4(t-\tau)}(1+\tau)^{-\frac{5}{2}}{\rm d}\tau\Big)^\frac{1}{2}\|b-u\|_{L_t^2(L^2)}\nonumber\\
\lesssim&\, (1+t)^{-\frac{5}{4}},
\end{align}
and
\begin{align}\label{NJKG3.163}
\|(\sqrt{2}\omega-\sqrt{3}\theta)(t)\|_{L^2}\lesssim&\, e^{-5t} \|\sqrt{2}\omega_0
-\sqrt{3}\theta_0\|_{L^2}+\int_0^t e^{-5(t-\tau)}\big(\|\nabla(u,b)\|_{L^2}
+\|\nabla f\|_{L_{x,v}^2}\big)  {\rm d}\tau\nonumber\\
&+\int_0^t e^{-5t}\big(\|b\cdot
 u\|_{L^2}+\|a\theta\|_{L^2}+\|\theta{\rm div} u\|_{L^2}+\|u\cdot\nabla\theta\|_{L^2} 
 \big){\rm d}\tau \nonumber\\
&+\int_0^t\big(  \|a\|_{L^6}\|u\|_{L^6}^2+\|u\cdot(b-u)\|_{L^2}  
+\|\rho\|_{L^\infty}\| \sqrt{2}\omega-\sqrt{3}\theta \|_{L^2} \big)
  {\rm d}\tau\nonumber\\
\lesssim&\, e^{-2t}+\int_0^t e^{-2(t-\tau)}(1+\tau)^{-\frac{5}{4}}{\rm d}\tau \nonumber\\
&+\Big(\int_0^t e^{-4(t-\tau)}(1+\tau)^{-\frac{5}{2}}{\rm d}\tau\Big)^\frac{1}{2}
\|\sqrt{2}\omega-\sqrt{3}\theta\|_{L_t^2(L^2)}\nonumber\\
\lesssim&\, (1+t)^{-\frac{5}{4}}.
\end{align}

We now proceed to derive the time-decay estimate  for the microscopic
 component $\{\mathbf{I}-\mathbf{P}\}f$ in \eqref{NJKd3}. To this end, we consider the equation 
 \eqref{Q2}$_4$, from which we obtain  
\begin{align}\label{NJKG3.164}
&\partial_t \{\mathbf{I}-\mathbf{P}\}f-\mathcal{L}{\{\mathbf{I}-\mathbf{P}\}}f
+v\cdot\nabla\{\mathbf{I}-\mathbf{P}\}f-\frac{1}{2}u\cdot v\mathbf{P}f+u\cdot\nabla_v\mathbf{P}f\nonumber\\
&\quad +\partial_t {\mathbf{P}f}+v\cdot\nabla {\mathbf{P}f}+(b-u)v\sqrt{M}
+\frac{1}{\sqrt{3}}(|v|^2-3)\sqrt{M}(\sqrt{2}\omega-\sqrt{3}\theta)\nonumber\\
&\qquad =  \frac{1}{2}u\cdot v\{\mathbf{I}-\mathbf{P}\}f
-u\cdot \nabla_v\{\mathbf{I}-\mathbf{P}\}f+\frac{\theta}{\sqrt{M}} \Delta_v(\sqrt{M}  \mathbf{P}f)\nonumber\\
& \qquad \ \ \, +\frac{\theta}{\sqrt{M}} \Delta_v(\sqrt{M} \{\mathbf{I}- \mathbf{P}\}f).
\end{align}
Taking the $L^2_{x,v}$ inner product of \eqref{NJKG3.164} and utilizing the
 property \eqref{G2.6} and Young's inequality, we  arrive at
\begin{align*} 
&\frac{1}{2}  \frac{\rm d}{{\rm d}t} \|\{\mathbf{I}-
\mathbf{P}\}f\|_{L_{x,v}^2}^2+\bar\lambda \|\{\mathbf{I}-\mathbf{P}\}f\|_{\nu}^2\nonumber\\
&\quad \lesssim \|\partial_t{\mathbf{P}f} \|_{L_{x,v}^2}^2
+\big(1+\|(u,\theta)\|_{H^2}^2\big)\|\nabla(a,b,\omega)\|_{L^2}^2+\|b-u\|_{L^2}^2
\nonumber\\
&\qquad+\|\sqrt{2}\omega-\sqrt{3}\theta\|_{L^2}^2 
+\eps_1\bar\lambda\|\{\mathbf{I}-\mathbf{P}\}f\|_{\nu}^2+\|u\|_{H^2}\|\{\mathbf{I}-\mathbf{P}\}f\|_{\nu}^2,
\end{align*}
which, together with the fact $ \|\{\mathbf{I}-\mathbf{P}\}f\|_{L_{x,v}^2} 
\lesssim\|\{\mathbf{I}-\mathbf{P}\}f\|_{\nu}^2$, leads to
\begin{align}\label{NJKG3.165}
  \|\{\mathbf{I}-\mathbf{P}\}f\|_{L_{x,v}^2}^2 
 \lesssim&\, \int_0^t e^{-2(t-\tau)}\big( \|\partial_t{\mathbf{P}f} 
 \|_{L_{x,v}^2}^2+\|\nabla(a,b,\omega)\|_{L^2}^2\nonumber\\
 & \quad + \|b-u\|_{L^2}^2
 +\|\sqrt{2}\omega-\sqrt{3}\theta\|_{L^2}^2  \big){\rm d}\tau %\nonumber\\
  +e^{-2\bar\lambda t}  \|\{\mathbf{I}-\mathbf{P}\}f_0\|_{L_{x,v}^2}^2\nonumber\\
 \lesssim&\, \int_0^t e^{-2(t-\tau)}(1+\tau)^{-\frac{5}{2}}{\rm d}\tau
 +e^{-2\bar\lambda t}+ \int_0^t e^{-2(t-\tau)}  \|\partial_t{\mathbf{P}f} \|_{L_{x,v}^2}^2{\rm d}\tau\nonumber\\
 \lesssim&\,(1+t)^{-\frac{5}{2}} + \int_0^t e^{-2(t-\tau)}  \|\partial_t{\mathbf{P}f} \|_{L_{x,v}^2}^2{\rm d}\tau.
\end{align}

Through a direct computation, it follows from \eqref{G3.21}$_1$–\eqref{G3.21}$_3$ that  
\begin{align}\label{NJKG3.166}
\|\partial_t{\mathbf{P}f} \|_{L_{x,v}^2}\lesssim&\, \|\partial_t a\|_{L^2}+ 
  \|\partial_t b\|_{L^2} +\|\partial_t \omega\|_{L^2}\nonumber\\
\lesssim\,&\|\nabla (a,b,\omega)\|_{L^2}+\|\nabla f\|_{L_{x,v}^2}+\|b-u\|_{L^2}\nonumber\\
& +\|\sqrt{2}\omega-\sqrt{3}\theta\|_{L^2}+\|au\|_{L^2}+\|u\cdot b\|_{L^2}\nonumber\\
\lesssim\,& (1+t)^{-\frac{5}{4}},
\end{align}

Substituting \eqref{NJKG3.166} into \eqref{NJKG3.165}, we obtain 
\begin{align}\label{NJKG3.167}
 \|\{\mathbf{I}-\mathbf{P}\}f\|_{L_{x,v}^2} \lesssim (1+t)^{-\frac{5}{4}}.   
\end{align}
By putting the estimates \eqref{NJKG3.162}, \eqref{NJKG3.163}, 
and \eqref{NJKG3.167} together, we obtain the desired  inequality \eqref{NJKd3}.
 Thus, the proof of Theorem \ref{T1.4} is completed.  
\end{proof}

\end{document}
